% PDT-Lean — blueprint content.
% Macros \lean{}, \leanok, \uses{} are provided by leanblueprint.
%   \lean{Decl}  links a statement to ONE Lean declaration `Decl` (one per node, so the
%                graph names every node cleanly; multiple decls per statement break naming)
%   \leanok      marks the statement as fully formalized (green in the graph)
%   \uses{...}   declares dependencies (the edges of the graph)
% `ass:posit` deliberately carries NO \leanok — the one identification the kernel does not
% evaluate (amber). The quantum statements are theorems about C and do not use it.

\chapter{Quantum kinematics from a complex line}

\begin{definition}[State space as a complex line]
  \label{def:qishilbert}
  \lean{PDT.QisHilbert}
  \leanok
  For a real vector space $H$, write $\mathtt{QisHilbert}\,H$ for the type of
  $\mathbb{R}$-linear equivalences $H \simeq \mathbb{C}$: the statement that $H$ carries a
  complex structure.
\end{definition}

\begin{assumption}[The one identification (not evaluated here)]
  \label{ass:posit}
  \uses{def:qishilbert}
  A physical quantum state space carries the complex structure of \cref{def:qishilbert}, and
  PDT identifies that $\mathbb{C}$ with the complex place of $Q$'s field --- the same $Q$ that
  fixes spacetime, so quantum mechanics and spacetime read off one arithmetic. The
  kinematics below are theorems about $\mathbb{C}$ and do not use this identification; whether the physical world's
  complex structure \emph{is} $Q$'s place is the single step \emph{not evaluated here}. (No
  \texttt{\textbackslash leanok}, by design: the kernel certifies the mathematics, not the
  identification.)
\end{assumption}

\begin{lemma}[Born form is the field norm]
  \label{lem:born_form}
  \lean{PDT.born_form}
  \leanok
  $(z\cdot z^{\dagger}).\mathrm{re} = \mathrm{normSq}\,z = |z|^2$, where $z^{\dagger}$ is the
  Galois conjugation $\mathtt{starRingEnd}\,\mathbb{C}$. The square is forced by
  $[\mathbb{C}:\mathbb{R}] = 2$.
\end{lemma}

\begin{lemma}[Hermitian = Galois-fixed = real (the pointer basis)]
  \label{lem:hermitian}
  \lean{PDT.hermitian_real}
  \leanok
  A Galois-fixed scalar is real: $z^{\dagger} = z \implies \mathrm{im}\,z = 0$.
\end{lemma}

\begin{lemma}[Unitary preserves the Born norm]
  \label{lem:unitary}
  \lean{PDT.unitary_normSq_one}
  \leanok
  $z\cdot z^{\dagger} = 1 \implies \mathrm{normSq}\,z = 1$.
\end{lemma}

\begin{theorem}[Born probabilities on $\mathbb{C}$]
  \label{thm:qm_from_q}
  \lean{PDT.QM_from_Q}
  \leanok
  \uses{lem:born_form, lem:hermitian, lem:unitary}
  The capstone, assembling the $\mathbb{C}$-lemmas above. For normalized two-outcome amplitudes $a,b\in\mathbb{C}$,
  the Born probabilities equal the conjugate-twisted trace form $|z|^2$, are
  non-negative, and sum to one --- a genuine probability rule, with \emph{no} independent
  Born postulate.
\end{theorem}

\chapter{The golden arithmetic bedrock}

\begin{proposition}[The ``23'' integer weld]
  \label{prop:disc_cubic}
  \lean{PDT.twentythree_weld}
  \leanok
  The integer identity $-\bigl(-4(-1)^3-27(-1)^2\bigr) = (3^2-1)+(4^2-1)$, both $=23$.
\end{proposition}

\begin{corollary}[Cubic discriminant and the gauge dimension (prose)]
  \uses{prop:disc_cubic}
  The left-hand side of \cref{prop:disc_cubic} is $-\operatorname{disc}(x^3-x-1)$
  (depressed-cubic formula $-4p^3-27q^2$, $p=q=-1$) and the right-hand side is
  $\dim\mathfrak{su}(3)+\dim\mathfrak{su}(4)$. Those two identifications are standard and are
  \emph{not} formalized ($\dim\mathfrak{su}(n)=n^2-1$ appears nowhere in the Lean development);
  only the integer identity of \cref{prop:disc_cubic} is kernel-checked. (No
  \texttt{\textbackslash leanok}: prose.)
\end{corollary}

\begin{proposition}[Quartic discriminant]
  \label{prop:disc_quartic}
  \lean{PDT.disc_quartic}
  \leanok
  $\operatorname{disc}(x^4-x-1) = -283$ (as the integer identity
  $-27(-1)^4+256(-1)^3 = -283$).
\end{proposition}

\begin{corollary}[Field discriminant of $\mathbb{Q}(Q)$ (prose)]
  \uses{prop:disc_quartic}
  Since $283$ is prime (\lstinline{PDT.prime_283}), $-283$ is squarefree, hence $\mathbb{Z}[Q]$
  is maximal and $-283$ is \emph{also} the field discriminant of $\mathbb{Q}(Q)$. This
  is standard; the chain squarefree $\Rightarrow$ maximal-order $\Rightarrow$ field-discriminant
  is a \emph{prose inference, not formalized}. (No \texttt{\textbackslash leanok}.)
\end{corollary}

\begin{proposition}[The compositum-norm numeral identity]
  \label{prop:norms}
  \lean{PDT.unit_norm_rhoQ}
  \leanok
  The numeral identity $1^4\cdot(-1)^3 = -1$. Its intended reading is the norm-tower product
  $N(\rho Q) = N(\rho)^4\cdot N(Q)^3 = -1$ (the unit norm forcing integer exponents), but the
  compositum $\mathbb{Q}(\rho,Q)$, the $[K:\mathbb{Q}]=12$ linear disjointness, and the
  norm-tower step are \emph{prose, not formalized} (see WO-4(ii)); only the numeral identity is
  kernel-checked. The genuine field norms $N(Q)=-1$ and $N(\rho)=1$ are stamped separately in
  \cref{prop:norm_q,prop:norm_r}; $N(2Q-1)=-23$ is the numeral identity \lstinline{PDT.norm_2Q_sub_one}
  (welding the ``$23$'' to $E_4$'s $2$-torsion --- also prose).
\end{proposition}

\chapter{The same arithmetic, through Mathlib's number-theory API}

The previous chapter states the norm and discriminant values as elementary integer
identities (closed by \texttt{norm\_num}). This chapter re-derives the \emph{same} facts
through Mathlib's genuine number-theory API --- \texttt{Algebra.norm}, \texttt{Algebra.discr}
and \texttt{Algebra.traceForm} over \texttt{AdjoinRoot}. This is the answer to ``the content
is in the names'': they are theorems about the real field-theoretic invariants, and each
rests on the polynomial being irreducible (so the \texttt{AdjoinRoot} quotient is a field).

\begin{proposition}[Cubic irreducible]
  \label{prop:irred_c}
  \lean{PDT.cubicQ_irreducible}
  \leanok
  $x^3-x-1$ is \texttt{Irreducible} over $\mathbb{Q}$, so $\mathbb{Q}[x]/(x^3-x-1)$ is a
  genuine number field.
\end{proposition}

\begin{proposition}[Quartic irreducible]
  \label{prop:irred_q}
  \lean{PDT.quarticQ_irreducible}
  \leanok
  $x^4-x-1$ is \texttt{Irreducible} over $\mathbb{Q}$, so $\mathbb{Q}[x]/(x^4-x-1)$ is a
  genuine number field.
\end{proposition}

\begin{proposition}[Norm of $Q$ via \texttt{Algebra.norm}]
  \label{prop:norm_q}
  \lean{PDT.norm_Q}
  \leanok
  \uses{prop:irred_q}
  $\operatorname{Algebra.norm}_{\mathbb{Q}}(Q) = -1$: $Q$ is a unit of norm $-1$, computed by
  the genuine \texttt{Algebra.norm} (agreeing with \cref{prop:norms}).
\end{proposition}

\begin{proposition}[Norm of $\rho$ via \texttt{Algebra.norm}]
  \label{prop:norm_r}
  \lean{PDT.norm_ρ}
  \leanok
  \uses{prop:irred_c}
  $\operatorname{Algebra.norm}_{\mathbb{Q}}(\rho) = 1$: $\rho$ is a unit of norm $+1$.
\end{proposition}

\begin{proposition}[Cubic discriminant via \texttt{Algebra.discr}]
  \label{prop:discr_c}
  \lean{PDT.discr_cubic}
  \leanok
  \uses{prop:irred_c}
  $\operatorname{Algebra.discr}_{\mathbb{Q}}$ of the cubic power basis $= -23$ --- the genuine
  discriminant of the basis (agreeing with \cref{prop:disc_cubic}).
\end{proposition}

\begin{proposition}[Quartic discriminant via \texttt{Algebra.discr}]
  \label{prop:discr_q}
  \lean{PDT.discr_quartic}
  \leanok
  \uses{prop:irred_q}
  $\operatorname{Algebra.discr}_{\mathbb{Q}}$ of the quartic power basis $= -283$ (agreeing
  with \cref{prop:disc_quartic}).
\end{proposition}

\begin{proposition}[The trace form \emph{is} \texttt{Algebra.traceForm}]
  \label{prop:traceform}
  \lean{PDT.traceForm_eq_M4}
  \leanok
  \uses{prop:irred_q, thm:sig_quartic}
  The Gram matrix $M_4$ behind the $(3,1)$ signature is \emph{exactly}
  $\operatorname{Algebra.traceForm}_{\mathbb{Q}}$ on the power basis of
  $\mathbb{Q}[x]/(x^4-x-1)$: $\operatorname{traceForm}(e_i,e_j) = (M_4)_{ij}$. Together with the
  kernel-level identity $M = M_4$ (\lstinline{PDT.M_eq_M4}, in \lstinline{PdtLinks}), the
  signature theorem is a statement about the genuine trace form, not a hand-built matrix.
\end{proposition}

\chapter{Trace-form signatures: 3-space and spacetime}

\begin{theorem}[Spacetime signature certificate (quartic)]
  \label{thm:sig_quartic}
  \lean{PDT.signature_3_1}
  \leanok
  The trace form $M$ of $\mathbb{Q}[x]/(x^4-x-1)$, in the power basis, is symmetric and there is a
  unimodular rational $P$ with $P^{\mathsf T} M P = \operatorname{diag}(4,4,-\tfrac94,\tfrac{283}{36})$
  (three positive diagonal entries, one negative) and $\det M = -283 < 0$. (A
  \emph{Sylvester congruence certificate over $\mathbb{Q}$}; the statement contains no
  signature vocabulary and no $\mathbb{R}$ base-change.)
\end{theorem}

\begin{corollary}[Signature $(3,1)$ (prose)]
  \uses{thm:sig_quartic}
  Base-changing to $\mathbb{R}$ and applying Sylvester's law of inertia, the certificate of
  \cref{thm:sig_quartic} gives signature $(3,1)$. This final step is standard and carried in
  prose, not formalized (see WO-4(i) for the route to a genuinely stamped
  $\operatorname{sigPos}/\operatorname{sigNeg}$ node). (No \texttt{\textbackslash leanok}.)
\end{corollary}

\begin{theorem}[$3$-space signature certificate (cubic)]
  \label{thm:sig_cubic}
  \lean{PDT.signature_2_1}
  \leanok
  The trace form $M_\rho$ of $\mathbb{Q}[x]/(x^3-x-1)$ is symmetric and congruent over
  $\mathbb{Q}$, via a unimodular $P_\rho$, to $\operatorname{diag}(3,2,-\tfrac{23}{6})$ --- two
  positive, one negative --- with $\det M_\rho = -23$. Moreover $M_\rho$ \emph{is} the genuine
  $\operatorname{Algebra.traceMatrix}$ Gram matrix of the cubic power basis
  (\lstinline{PDT.Mρ_eq_traceMatrix_fin3}, in \lstinline{PdtLinks}). (A Sylvester congruence
  certificate over $\mathbb{Q}$.)
\end{theorem}

\begin{corollary}[Signature $(2,1)$ (prose)]
  \uses{thm:sig_cubic}
  As in the quartic case, base-change to $\mathbb{R}$ and Sylvester's law give signature $(2,1)$;
  this final step is standard and carried in prose, not formalized. (No
  \texttt{\textbackslash leanok}.)
\end{corollary}

\chapter{The Bell gap: classical $2$ versus quantum $2\sqrt2$}

\begin{theorem}[Classical CHSH bound]
  \label{thm:chsh_classical}
  \lean{PDT.chsh_comm_le_two}
  \leanok
  In a commutative ordered $*$-ring, any CHSH tuple satisfies
  $A_0B_0 + A_0B_1 + A_1B_0 - A_1B_1 \le 2$ (via Mathlib's
  \lstinline{CHSH_inequality_of_comm}). Commuting observables stop at $2$.
\end{theorem}

\begin{theorem}[Tsirelson bound, saturated]
  \label{thm:tsirelson}
  \lean{PDT.chsh_saturates}
  \leanok
  The Tsirelson bound $2\sqrt2$ (Mathlib's \lstinline{tsirelson_inequality}) is attained:
  an explicit real-Pauli quadruple in $M_2(\mathbb{R})\otimes M_2(\mathbb{R})$ is a genuine
  CHSH tuple whose operator reaches eigenvalue $2\sqrt2$ on the Bell direction. The gap
  $2 < 2\sqrt2$ (\lstinline{PDT.bell_gap}) is strict. (This node is independent of
  \cref{thm:chsh_classical}: \lstinline{chsh_saturates} makes no use of the classical bound,
  so there is no formal dependency edge.)
\end{theorem}


\chapter{Two forms and one rotation on the complex place}

\begin{theorem}[The dagger identity]
  \label{thm:dagger_identity}
  \lean{PDT.born_eq_trace_comp_conj}
  \leanok
  $G_{\mathrm{Born}} = G_{\mathrm{trace}} \circ \sigma$: the Born form is the trace form
  precomposed with conjugation.
\end{theorem}

\begin{theorem}[$J$ is a Born isometry and a trace anti-isometry]
  \label{thm:j_isometries}
  \lean{PDT.J_born_isometry}
  \leanok
  \uses{thm:dagger_identity}
  The same $J$ (multiplication by $i$) satisfies $J^{\top} G_{\mathrm{Born}} J = G_{\mathrm{Born}}$
  and $J^{\top} G_{\mathrm{tr}} J = -G_{\mathrm{tr}}$ (companion: \texttt{J\_trace\_antiisometry}).
\end{theorem}

\begin{theorem}[The symplectic companion]
  \label{thm:omega_symplectic}
  \lean{PDT.omega_born_nondegenerate}
  \leanok
  \uses{thm:j_isometries}
  $\omega = J^{\top} G_{\mathrm{Born}}$ is alternating and nondegenerate ($\det = 1$) — a symplectic
  form — while the trace companion $J^{\top} G_{\mathrm{tr}}$ is symmetric (companions:
  \texttt{omegaB\_eq}, \texttt{omega\_born\_skew}, \texttt{omega\_trace\_symm}).
\end{theorem}

\chapter{Anchored stabilizer dimensions}

\begin{theorem}[Surjectivity of the anchored action]
  \label{thm:slaction_surj}
  \lean{PdtStabilizer.slAction_surjective}
  \leanok
  For $N \ge 2$ and $v \ne 0$, the map $X \mapsto Xv$ from trace-zero matrices onto
  $\mathbb{Q}^N$ is surjective, by an explicit solver.
\end{theorem}

\begin{theorem}[The stabilizer has dimension $N^2-N-1$]
  \label{thm:finrank_stab}
  \lean{PdtStabilizer.finrank_stabilizer}
  \leanok
  \uses{thm:slaction_surj}
  The sl-stabilizer of $v \ne 0$ has finrank $N^2-N-1$ (rank–nullity); instances:
  $11$ at $N{=}4$, $19$ at $N{=}5$, $209$ at $N{=}15$.
\end{theorem}

\begin{theorem}[Corank one: the trace character]
  \label{thm:corank_one}
  \lean{PdtStabilizer.stabilizer_corank_one}
  \leanok
  \uses{thm:finrank_stab}
  The gl-stabilizer has finrank $N^2-N$; the corank-1 gap to the sl-stabilizer is the trace
  character, exhibited by an explicit trace-one witness.
\end{theorem}

\begin{theorem}[Orbit--stabilizer decomposition]
  \label{thm:orbit_stab}
  \lean{PdtStabilizer.finrank_sl_eq_add}
  \leanok
  \uses{thm:finrank_stab}
  $\dim \mathfrak{sl}(N) = N + (N^2-N-1)$.
\end{theorem}

\chapter{The Pisot boundary}

Both statements are proved by a single Vieta reduction on the conjugate pair, from the
defining equation alone --- no root-finding and no \lstinline{Polynomial} API. A non-real
root is characterised only by $z^3 = z+1$, $\operatorname{im} z \neq 0$ (resp.\ $w^4 = w+1$),
so each theorem quantifies over all such roots.

\begin{theorem}[The cubic's non-real roots lie inside the unit circle]
  \label{thm:cubic_lt_one}
  \lean{PDT.cubic_conj_norm_lt_one}
  \leanok
  Every non-real root $z$ of $x^3-x-1$ has $\|z\| < 1$. Equivalently: $\rho$ is a Pisot
  number.
\end{theorem}

\begin{theorem}[The quartic's non-real roots lie outside the unit circle]
  \label{thm:quartic_gt_one}
  \lean{PDT.quartic_conj_norm_gt_one}
  \leanok
  Every non-real root $w$ of $x^4-x-1$ has $\|w\| > 1$. Equivalently: $Q$ is not a Pisot
  number.
\end{theorem}

\begin{theorem}[The Pisot-boundary dichotomy]
  \label{thm:pisot_dichotomy}
  \lean{PDT.pisot_boundary_dichotomy}
  \leanok
  \uses{thm:cubic_lt_one, thm:quartic_gt_one}
  The two statements packaged as one: $\|z\| < 1$ and $\|w\| > 1$. The two zero-parameter
  polynomials sit on opposite sides of the unit circle. This upgrades the computed moduli
  $0.869$ and $1.063$ to a theorem.
\end{theorem}

\begin{assumption}[Reading the dichotomy dynamically (not evaluated here)]
  \label{ass:settle_spiral}
  \uses{thm:pisot_dichotomy}
  Reading $\|z\| < 1$ as a state that \emph{settles} and $\|w\| > 1$ as one that
  \emph{spirals} is an \emph{identification}. What the kernel proves is the two modulus
  inequalities. (No \texttt{\textbackslash leanok}, by design.)
\end{assumption}

\chapter{The $\beta$-clock word}

For a real slope $\beta$, the tick word is $m(n) = \lfloor (n+1)\beta\rfloor - \lfloor
n\beta\rfloor$ (\lstinline{PDT.m}). Irrationality is never assumed globally: it is carried as
an explicit hypothesis on the one statement that needs it.

\begin{theorem}[Two-valued]
  \label{thm:clock_two_valued}
  \lean{PDT.m_mem}
  \leanok
  For $0 < \beta < 1$, every $m(n)$ is $0$ or $1$: the word is over a binary alphabet.
\end{theorem}

\begin{theorem}[Density exactly $\beta$]
  \label{thm:clock_density}
  \lean{PDT.density_tendsto}
  \leanok
  $\bigl(\sum_{k<N} m(k)\bigr)/N \to \beta$, by exact telescoping of the partial sums to
  $\lfloor N\beta\rfloor$ and a squeeze.
\end{theorem}

\begin{theorem}[Balanced]
  \label{thm:clock_balanced}
  \lean{PDT.balanced}
  \leanok
  \uses{thm:clock_two_valued}
  Any two windows of the same length $L$ have $1$-counts differing by at most $1$ --- the
  Sturmian balance property; both counts lie in $\{\lfloor L\beta\rfloor, \lfloor
  L\beta\rfloor + 1\}$.
\end{theorem}

\begin{theorem}[Aperiodic, for irrational slope]
  \label{thm:clock_aperiodic}
  \lean{PDT.aperiodic}
  \leanok
  For irrational $\beta$ --- carried as the explicit hypothesis
  \lstinline{Irrational \beta} --- there is no $p > 0$ with $m(n+p) = m(n)$ for all $n$. A
  period would force $M\cdot\operatorname{fract}(p\beta) < 1$ for every $M$, contradicting the
  Archimedean property. Aperiodic but not random.
\end{theorem}

\chapter{Contraction and separation on the complex place of $\mathbb{Q}(\rho)$}

Throughout: $r$ is the real root of $x^3-x-1$ and $z$ a non-real root of the same cubic,
each characterised only by its defining equation, so every statement quantifies over all
such $r,z$. Write $\sigma_1$ for the real and $\sigma_2$ for the complex embedding of
$\mathbb{Q}(\rho)$, and identify an element of $\mathbb{Z}[\rho]$ with its coordinate
triple in the power basis $(1,\rho,\rho^2)$.

\begin{theorem}[Exact contraction rate]
  \label{thm:tick_contraction}
  \lean{PDT.tick_contraction}
  \leanok
  $r\cdot\|z\|^2 = 1$: multiplication by $z$ scales area on the complex place by exactly
  $1/r$. This is the norm relation $N(\rho)=1$ read at the two archimedean places --- an
  identity of the field, not a numerical estimate.
\end{theorem}

\begin{theorem}[Rate form: $\|z\| = r^{-1/2}$]
  \label{thm:norm_z_rpow}
  \lean{PDT.norm_z_eq_rpow}
  \leanok
  \uses{thm:tick_contraction}
  Lengths on the complex place contract by exactly $r^{-1/2}$ under multiplication by $z$;
  in particular $\|z\| < 1$ (\lstinline{PDT.norm_z_lt_one}), since $r > 1$.
\end{theorem}

\begin{theorem}[Integrality floor]
  \label{thm:meyer_separation}
  \lean{PDT.meyer_separation}
  \leanok
  \uses{prop:irred_c}
  For every nonzero $x \in \mathbb{Z}[\rho]$, $\;1 \le |\sigma_1(x)|\cdot\|\sigma_2(x)\|^2$.
  The product of the three embeddings of $x$ is the integer $N(x)$
  (\lstinline{PDT.normForm_factor}), and $N(x) \neq 0$ because $x^3-x-1$ is irreducible
  (\cref{prop:irred_c}); so the left-hand side is $|N(x)| \ge 1$.
\end{theorem}

\begin{theorem}[Radius form of the floor]
  \label{thm:meyer_rpow}
  \lean{PDT.meyer_separation_rpow}
  \leanok
  \uses{thm:meyer_separation}
  $\|\sigma_2(x)\| \ge |\sigma_1(x)|^{-1/2}$ for every nonzero $x \in \mathbb{Z}[\rho]$.
\end{theorem}

\begin{theorem}[Windowed uniform discreteness]
  \label{thm:windowed_separation}
  \lean{PDT.windowed_separation}
  \leanok
  \uses{thm:meyer_rpow}
  Two distinct points of $\mathbb{Z}[\rho]$ whose real embeddings differ by at most $H$ have
  complex embeddings at distance at least $H^{-1/2}$.
\end{theorem}

\begin{theorem}[Windowed permanence]
  \label{thm:certification}
  \lean{PDT.certification_soundness}
  \leanok
  \uses{thm:windowed_separation, thm:norm_z_rpow}
  If $\|w - \sigma_2(\ell)\| < \tfrac12 H^{-1/2}$ then for every $n$ the point $z^n w$ is
  strictly closer to $z^n\sigma_2(\ell) = \sigma_2(\rho^n\ell)$ than to any other
  $\sigma_2(\ell')$ whose real embedding lies within $H$ of that of $\rho^n\ell$
  (multiplication by $z$ permutes $\sigma_2(\mathbb{Z}[\rho])$,
  \lstinline{PDT.lattice_tick_invariant}). The window hypothesis is essential: the
  unwindowed image $\sigma_2(\mathbb{Z}[\rho])$ is dense in $\mathbb{C}$, so no unwindowed
  statement of this kind is true and none is claimed.
\end{theorem}

\begin{assumption}[The dynamical reading (not evaluated here)]
  \label{ass:record}
  \uses{thm:certification}
  Reading multiplication by $z$ as one step of a dynamics, $\sigma_2(\mathbb{Z}[\rho])$ as a
  record lattice, and \cref{thm:certification} as the permanence of a recorded outcome is an
  \emph{identification}. The theorems above are statements about two embeddings of
  $\mathbb{Q}(\rho)$ and the metric geometry of $\mathbb{C}$; the physical reading is stated
  outside the kernel. (No \texttt{\textbackslash leanok}, by design.)
\end{assumption}
